% Название разделов -- все прописные
\section{ТЕОРЕТИЧЕСКИЕ ОСНОВЫ И АНАЛИЗ ПРЕДМЕТНОЙ ОБЛАСТИ}

Разработка системы индексации ончейн-данных требует рассмотрения как
специфики блокчейн-сетей, так и особенностей построения прикладных
информационных систем, работающих с большими потоками событий. В отличие от
традиционных транзакционных систем, блокчейн Bitcoin выступает одновременно
распределенным журналом, средой консенсуса и источником данных для внешних
программных сервисов. Это накладывает ограничения на способы извлечения
данных, их интерпретацию, хранение и последующее предоставление через
прикладные интерфейсы \cite{narayanan2016bitcoin, antonopoulos2017mastering}.

Для обоснования проектных решений по теме ВКР необходимо определить сущность
ончейн-данных, выявить особенности блокчейна Bitcoin как источника
информации, проанализировать ограничения прямой работы прикладных систем с
узлом сети, а также сопоставить существующие подходы к индексации
блокчейн-данных. Результат такого анализа позволяет не только сформулировать
задачу разработки системы Bitcoin Blockchain Indexer, но и
доказательно обосновать необходимость выделения специализированного слоя
индексации и хранения данных.

\subsection{Анализ существующих подходов к индексации ончейн-данных}

Задача предоставления прикладного доступа к данным Bitcoin может решаться
разными способами: прямым использованием RPC узла, развертыванием готового
индексатора, применением аналитических платформ или выгрузкой данных в
пакетном ETL-формате. Каждый подход закрывает часть требований, однако не все
из них одинаково подходят для системы, в которой необходимы управляемые
задания индексирования, REST API, PostgreSQL-хранилище, контроль reorg,
mempool и эксплуатационная наблюдаемость.

Bitcoin Core предоставляет канонический источник данных и RPC-интерфейс, но
его назначение состоит прежде всего в поддержании узла сети, а не в
обслуживании прикладных выборок по адресам и агрегатам \cite{btc_core_rpc}.
Индексаторы семейства electrs ориентированы на быстрые запросы кошельков и
используют собственное хранилище индексов; проект Esplora дополняет такой
подход web-интерфейсом и HTTP API для обозревателя блоков
\cite{electrs_repo, esplora_repo}. Аналитическая платформа BlockSci
предназначена для исследовательского анализа блокчейн-данных и использует
аналитическую модель хранения, но не является транзакционным backend-сервисом
для оперативной выдачи данных внешним приложениям \cite{blocksci_paper}.
Проект Bitcoin ETL решает задачу пакетной выгрузки блоков и транзакций,
включая сценарии потоковой передачи, однако его модель ближе к подготовке
датасетов, чем к сопровождению управляемого индексатора с API
\cite{bitcoin_etl_repo}. Публичные explorer API удобны для быстрых интеграций,
но не дают полного контроля над источником данных, политикой хранения,
ограничениями доступа и воспроизводимостью результатов.

\begin{footnotesize}
\setlength{\tabcolsep}{3pt}
\begin{longtable}{|p{0.17\textwidth}|p{0.20\textwidth}|p{0.20\textwidth}|p{0.11\textwidth}|p{0.17\textwidth}|}
    \caption{Сравнение подходов к получению и индексации данных Bitcoin}\label{tab:indexer-approaches}\\ \hline
    Подход & Сильные стороны & Ограничения & Контроль & Вывод \\ \hline
    \endfirsthead
    \caption*{Продолжение таблицы \ref{tab:indexer-approaches}}\\ \hline
    Подход & Сильные стороны & Ограничения & Контроль & Вывод \\ \hline
    \endhead
    \hline
    \endfoot
    \hline
    \endlastfoot
    Bitcoin Core RPC &
    Канонический источник блоков и транзакций &
    Нет готовых агрегатов и управления заданиями &
    высокий &
    только как источник данных \\ \hline
    electrs / Esplora &
    Быстрые адресные запросы и explorer API &
    Иная модель хранения, фокус на обозревателе &
    средний &
    полезен как аналог, но не закрывает ТЗ \\ \hline
    BlockSci &
    Аналитическая модель для исследований &
    Не является оперативным REST backend &
    высокий &
    подходит для анализа, не для эксплуатации \\ \hline
    Bitcoin ETL &
    Пакетная выгрузка и подготовка датасетов &
    Нет постоянного API и жизненного цикла заданий &
    средний &
    применим для выгрузок, но не как основной backend \\ \hline
    Публичные explorer API &
    Быстрый старт без собственного узла &
    Зависимость от внешнего сервиса и лимитов &
    низкий &
    не подходит для ВКР \\ \hline
\end{longtable}
\end{footnotesize}

Проведенное сравнение показывает, что готовые решения полезны как ориентиры,
но целевая система требует собственного слоя индексации. Выбор модульного
stateless-монолита позволяет сохранить простоту развертывания MVP, выделить
ответственность модулей внутри одного backend-приложения и хранить состояние в
PostgreSQL, а не в памяти процесса. Такой подход согласуется с требованиями к
воспроизводимости, тестируемости и дальнейшему поэтапному развитию системы.
